\chapter*{Úvod}
Cílem této práce je navrhnout a implementovat herní mechaniky pro postavy videohry žánru 2D plošinovky. Práce se zaměřuje na systémy ovládání hráče a návrh logiky nepřítele. Mechaniky hráče nezahrnují pouze pohyb do stran, ale také systém skoku a střelby. Zvláštní důraz je kladen na implementaci skoku, jehož správné nastavení zásadně ovlivňuje herní pocit a odlišuje kvalitní 2D plošinovky od méně povedených titulů.

Součástí práce je také implementace fyzikálního systému. Herní engine Unity sice nabízí vestavěné nástroje pro práci s fyzikou, ty však nemusí vždy poskytovat požadovanou míru kontroly nad chováním postavy. Z tohoto důvodu jsou tyto nástroje využity především jako základ, zatímco pokročilejší herní mechaniky jsou řešeny vlastním kódem.

S programováním herních mechanik úzce souvisí také implementace animací. Po vytvoření grafických podkladů je nutné jednotlivé animace správně integrovat do enginu a propojit je s řídicí logikou tak, aby odpovídaly aktuálnímu stavu postavy.

Významnou část projektu tvoří rovněž systém střelby založen na automatickém míření a výstřelu při inputu hráče. Tento systém přináší řadu technických výzev, například správné určení směru projektilu, úpravu jeho rychlosti nebo zajištění jeho odstranění po určité době z důvodu optimalizace výkonu hry.
